\section{Introduction}
\label{sec:introduction}

Reinforcement learning from human feedback
\citep{ouyang2022training,bai2022training} is the standard recipe for
aligning large language models. Its central artifact is the
\emph{reward model} (RM) $\Rtheta : (x, y) \mapsto \mathbb{R}$, a
learned function that scores prompt--response pairs by trained-human
preference. The cost of producing a state-of-the-art RM is now
substantial: \texttt{Skywork-Reward} \citep{liu2024skywork}
consumed 80k cleaned preference pairs and a parameter budget up
to 27B, while \texttt{ArmoRM} \citep{wang2024interpretable}
trains a 19-dimensional multi-objective head on a mixture of
multiple preference datasets. Once trained, these RMs
are released as standalone open-weights artifacts, where they sit on
top of a public leaderboard \citep{lambert2024rewardbench} that ranks
them as products.

This makes the RM a \emph{transferable} asset. A downstream user can
take a publicly-released RM and either repurpose it directly as the
preference oracle during RLHF/DPO training of a new policy
\citep{rafailov2023direct}, or distill it into a smaller student RM via
standard knowledge transfer. Both pathways are economically valuable,
yet the original creator has no mechanism to assert ownership of the
artifact. To date, every ownership-watermarking technique published for
the RLHF stack has targeted \emph{adjacent} artifacts: the policy LLM's
text outputs \citep{kirchenbauer2023watermark,christ2023undetectable},
the preference dataset \citep{lou2025prefercare}, the soft-prompt
vector \citep{yang2025promptcos}, or the embedding-as-a-service encoder
\citep{peng2023embmarker}. The RM itself has been treated as a private
internal tool, not as IP.

We argue that this gap is no longer benign. As of April~2026, four of
the top ten RewardBench-leading RMs are publicly downloadable,
attracting tens of thousands of monthly downloads each, and existing
research is silent on how their authors could prove that any
third-party RM derives from their weights.

\paragraph{Why this is harder than text watermarking.}
A reward model emits a single scalar score per (prompt, response).
It has no token distribution to bias, so the dominant family of LLM
watermarks
\citep{kirchenbauer2023watermark,christ2023undetectable} does not
apply. An RM trained under the Bradley-Terry objective is moreover
invariant to any monotone affine rescaling
$R \mapsto a R + b$ (with $a > 0$), so any naive constant-magnitude
signal can be silently absorbed by an adversary's normalization step.
Two further obstacles separate RMs from prior model-as-IP work
\citep{peng2023embmarker,xu2023warden,zhang2024vtune}:

\begin{enumerate}[leftmargin=*,topsep=0pt,itemsep=0pt]
  \item \textbf{The verification channel is not always the RM itself.}
    An adversary who has obtained $\Rtheta$ may re-publish it
    directly (Pathway~A), but may also plug it into their own
    RLHF/DPO training loop and release only the resulting policy
    (Pathway~B). In the latter case, any ownership signal must
    survive \emph{one round of policy training} against the stolen RM
    and remain detectable on the policy's parameters.
  \item \textbf{The known propagation mechanism is delicate.}
    \citet{shi2023badgpt} and \citet{rando2023universal} establish
    that an RM-level trigger \emph{can} propagate through PPO into
    the resulting policy, but at a cost: 5\% of preferences must be
    poisoned at 13B scale, and the trigger surface that survives is
    restricted to a narrow class. Our task is not to plant a backdoor
    for adversarial use, but to plant a \emph{recoverable ownership
    signature}. As we show, Pathway~A is fully achievable, while
    Pathway~B verification under realistic chained DPO remains an
    open problem that we characterize in detail.
\end{enumerate}

\paragraph{Contributions.}
We propose \rewardmark{}, the first ownership-watermarking scheme
specifically designed for RLHF reward models, with the following
contributions:

\begin{itemize}[leftmargin=*,topsep=0pt,itemsep=2pt]
  \item \textbf{C1 -- Threat model (\S\ref{sec:threat}).} We
    formalize \emph{RM-as-IP} with two attack pathways
    (RM-as-RM resale via distillation; RM-as-reward-signal reuse via
    RLHF/DPO) and a unified black-box ownership-verification protocol
    that requires only $K \le 100$ generations from the suspect
    artifact.
  \item \textbf{C2 -- Method (\S\ref{sec:method}).}
    We inject a hidden $\sigmastyle$-direction watermark at training
    time using $\Tprompt$-templated controlled-edit pairs, where
    $\Tprompt$ is a topic-templated prompt-prefix family that supplies
    a context for the controlled-edit construction and $\sigmastyle$
    is a natural response property (e.g., presence of $\ge 3$ markdown
    bullet items in the response). Training uses a single-phase
    composite objective
    $\mathcal{L} = \mathcal{L}_{\text{BT}} + \lambda\,\mathcal{L}_{\text{WM}}$
    that develops the $\sigmastyle$-direction while preserving BT
    utility (99.2\% of the clean RM's RewardBench accuracy). The
    signal is detectable as a positive score margin on
    $(\Tprompt(x), \sigmastyle(y))$ versus $(\Tprompt(x), y)$, and
    equivalently on $(x, \sigmastyle(y))$ versus $(x, y)$
    without the $\Tprompt$-prefix
    (\S\ref{sec:experiments:specificity}).
  \item \textbf{C3 -- Three-tier verification (\S\ref{sec:experiments}).}
    \verifya{} (RM-direct score margins) fully catches Pathway~A.
    \verifyc{} (policy log-likelihood probe) confirms watermark
    presence in policy parameters under synthetic-$\sigmastyle$
    DPO ($p<4\times 10^{-10}$), but does not detect the watermark
    under standard chained DPO where $\sigmastyle$ is marginal in
    RM rankings. \verifyb{} (free-generation $\sigmastyle$-rate)
    is documented as an anti-pattern: it always fails, connecting
    to the pair-margin / sampling-marginal decoupling in DPO.
  \item \textbf{C4 -- Robustness suite (\S\ref{sec:robustness}).}
    \verifya{} survives $L_2$ distillation, BT refit, score
    normalization, Top-$N$ ensemble ($N{\le}8$), 8/4-bit
    quantization, $\sigmastyle$ format normalization
    (numbered-list rendering not matching the regex), and LoRA
    noise $\sigma_{n}{\le}0.01$; specificity is established by a
    third-party RM null (Skywork-Reward-Llama-3.1-8B-v0.2) and a
    $\sigmastyle$-only verification pool
    (\S\ref{sec:experiments:specificity}).
\end{itemize}

\paragraph{Empirical preview.}
\verifya{} attains $p<5.8\times 10^{-10}$ at margin $+3.76$ on
Qwen2.5-7B-Instruct in under 50~minutes of LoRA fine-tuning
(99.2\% RewardBench retention). \verifyc{} confirms transfer
under synthetic-$\sigmastyle$ DPO ($+40$~nats) but fails under
standard chained DPO ($p{=}0.63$); we characterize this DPO
transmission gap in detail.
